\subsection{Overview}

Dissolved reactive phosphorus (DRP) exported through subsurface (tile) drains
is difficult to intercept because it moves in solution rather than attached to
sediment; reactive filter media that adsorb or precipitate phosphate offer an
edge-of-field means of retaining it before drainage water reaches surface waters.
The evidence base assembled here spans batch equilibrium tests, saturated column
and flow-through experiments, and full-scale phosphorus (P) removal structures,
and it converges on a two-family conceptual model. Iron- (Fe-)based media
(iron-oxide-coated sand, iron oxyhydroxides, zero-valent iron, iron filings,
and steel by-products) retain phosphate primarily by strong inner-sphere
adsorption, whereas calcium- (Ca-)rich media (blast-furnace and steel slag,
limestone, lime-amended pellets, seashells) retain it primarily by
Ca-phosphate precipitation driven by dissolved $\mathrm{Ca^{2+}}$ and elevated pH.
Across systematic comparisons, sorption is identified as the preferred DRP
retention mechanism at the edge of tile-drained catchments, and Fe-based
materials are generally described as more suitable than Ca-based materials
because they retain phosphate more stably and over a wider pH range
\parencite{Mendes2022,Li2020}. Headline findings are that laboratory batch and
column tests routinely achieve $\geq 90$--$99\%$ phosphate removal for the better
media, that industrial by-products adsorb roughly an order of magnitude more
phosphate than natural minerals \parencite{Sellner2019b}, and that field-scale
structures treating real drainage flow have removed on the order of
27--55\% of the dissolved P load \parencite{Penn2020,Penn2024,Shedekar2020}.
Field performance is consistently lower and more variable than laboratory
performance, with hydraulic retention time, pH, and cumulative P loading
identified as the controlling factors.

\subsection{Classes of Reactive Media and Phosphate Sorption Capacity}

\subsubsection{Iron-based media}

Iron-based materials are the best-performing and most consistently favoured
class. In replicated laboratory comparisons of four iron-based filter materials,
each removed at least 95\% of the soluble phosphorus originally present in
batch tests spiked with 1, 10, and 100~ppm $\mathrm{PO_4^{3-}}$-P; at 1000~ppm the
ranking of removal effectiveness was iron oxyhydroxide (IOH) $>$ sand--iron
mixture (SMI) $>$ zero-valent iron (ZVI) $>$ powdered iron composite (PIC), and
IOH was judged to hold a slight edge for agricultural drainage-water treatment
\parencite{Allred2014}. Earlier batch work showed at least 94\% phosphate removal
for all ZVI and SMI materials tested, and column tests confirmed substantial
phosphate removal by SMI3 and ZVI1 even at moderate-to-high concentrations and
very short exposure times \parencite{Allred2012}. Synthetic goethite iron
oxyhydroxide (SG-IOH) almost totally removed phosphate ($>98\%$) from drainage
waters spiked with an additional 1 or 10~ppm phosphate-P in saturated columns,
with the original formulation (SG-IOH-O) outperforming the regenerated one
\parencite{Allred2017}. Among natural and industrial adsorbents compared at
identical conditions, iron filings reached the highest capacity of 4.95~mg/g,
and steel slag 1.68~mg/g, versus only 0.10--0.13~mg/g for three natural minerals
at $20^{\circ}$C---an approximately order-of-magnitude advantage for the
industrial by-products \parencite{Sellner2019b}. In long-term column experiments,
the Fe-based CFH material exhibited the highest P retention (retaining 99\% of
high inlet P and 98\% of low inlet P) and was the most robust under variable flow
and loading, whereas calcined diatomite earth (CDE) performed moderately but
consistently \parencite{Pugliese2023}. Granular iron-oxide media (GNEM) removed
soluble reactive phosphorus (SRP) below detection limits far longer than steel
by-product media in continuous columns and reached a final saturated adsorption
capacity of 1.176~mg~P/g \parencite{Hauda2020}. Reviews reinforce this ordering:
Fe-based filter materials generally achieved DRP retention $\geq 95\%$ in batch
tests and $>90\%$ in most column tests, showed negligible desorption, and formed
strong, stable bonds \parencite{Mendes2022}. Iron-activated alumina (Alcan),
marketed for arsenic removal, achieved up to 94.4\% phosphate removal in batch
tests with a maximum capacity of 2.608~mg/g and was proposed as a novel DRP
adsorbent for subsurface drainage \parencite{Kumari2025}.

\subsubsection{Calcium-rich slag and lime-based media}

Calcium-rich media are widely available and effective, but their performance is
governed by the finite supply of reactive $\mathrm{Ca^{2+}}$ and buffering
capacity. Granulated blast-furnace slag (GBFS) and cement-kiln-dust (CKD)/silica
sand blends exhibited $>98\%$ $\mathrm{PO_4^{3-}}$ removal across batch experiments
spanning 3--68~mg/L initial concentration \parencite{Agrawal2011}. Furnace slag
media likewise removed all available P at initial concentrations of 1--5~mg~P/L,
with the fine particles removing all available P up to 100~mg~P/L and following
the sorption order FS\_Fine $>$ FS\_SM $>$ FS\_LG at 600~mg~P/L; in columns all
furnace slag treatments outperformed sandy-loam-coated pea gravel, with average
P removal of 50\% (FS\_SM), 22\% (FS\_LG) and 27\% (FS\_C) \parencite{Sanford2016}.
Steel slag reached 1.68~mg/g in batch tests, and mixtures of small steel chips
with limestone yielded 0.95, 1.42, and 2.28~mg/g at 10\%, 20\%, and 40\% steel
chip content, respectively \parencite{Sellner2019b}. Fly-ash ceramic pellets
amended with lime and bentonite clay were engineered to raise phosphate
adsorption while limiting heavy-metal leaching: a lime-free pellet reduced
solution P by 80\% within 30~s but leached excessive Cr, whereas the highest-lime
pellet had the greatest porosity (14\%) and saturated capacity (1.98~mg~P/g) and
the lowest Cr leaching (30~$\mu$g/L) \parencite{Li2017}. Reviews report that
Ca-based media commonly achieved 100\% DRP retention at low DRP concentrations but
declined at higher concentrations, and that some Ca-based materials (furnace slag,
Leca) showed lower column efficiencies deemed unsuitable for the field
\parencite{Mendes2022}; Filtralite-P and seashells performed well at high inlet P
but relatively poorly at low inlet P, and limestone showed the lowest retention
and highest sensitivity to loading \parencite{Pugliese2023}.

\subsubsection{Natural minerals and other engineered media}

Natural minerals and unmodified sorbents span a wide, generally lower,
performance range. Zeolite and silica sand gave statistically similar, modest
$\mathrm{PO_4^{3-}}$ removal of 21--58\%, and coconut-shell activated carbon
(CS-AC) removed 70--79\% \parencite{Agrawal2011}. Modified clinoptilolite zeolite
improved removal by a factor of 13--18 relative to raw samples, and in
drainage-model experiments zeolite placed around the drains achieved the highest
removals of the tested configurations (29.28\% phosphate, alongside 59.72\%
nitrate and 77.47\% salinity) \parencite{Sadeghi2021}. Comprehensive review of
inorganic adsorbents found that phosphate adsorption capacity varies widely among
classes and is strongly affected by pH, ionic strength, and coexisting ions; some
metal-loaded resins reached 2000--3000~mM/kg at near-neutral pH and generally
outperformed plain metal oxyhydroxides and ZVI, while ferrihydrite,
lepidocrocite, selected Zr/La materials, layered double hydroxides (LDHs), and
hybrid anion-exchange resins retained relatively strong performance near neutral
pH \parencite{Li2020}.

\subsection{Sorption and Precipitation Mechanisms}

The two media families retain phosphate by mechanistically distinct routes.
For Fe-based media, retention is dominated by inner-sphere adsorption that forms
strong, stable bonds and shows negligible desorption \parencite{Mendes2022}.
Isotherm and thermodynamic analyses support chemical binding: best-fit isotherms
were Dubinin--Radushkevich for GBFS and zeolite, Freundlich for silica sand,
Frumkin for CS-AC, and Langmuir for both CKD/sand blends, with mean free-energy
values indicating ion-exchange/chemical sorption for GBFS and the 10:90\% CKD/sand
blend but physical sorption for zeolite, silica sand, and CS-AC
\parencite{Agrawal2011}. For the CKD/sand blends the separation factor satisfied
$0 < R_L < 1$ and decreased with rising initial concentration, indicating more
favourable sorption at higher concentrations \parencite{Agrawal2011}. The Langmuir
model generally fit steel and mineral media better than Freundlich, and all
tested materials showed increasingly negative Gibbs free energy and positive
enthalpy with rising temperature, indicating endothermic adsorption
\parencite{Sellner2019b}. Alcan followed a Freundlich isotherm ($R^2=0.991$,
multilayer adsorption) and pseudo-second-order kinetics ($R^2=0.99$), with a
calculated $q_e$ close to the experimental value supporting chemisorption as the
dominant mechanism \parencite{Kumari2025}. For lime-amended pellets the
Langmuir--Freundlich model fit best, with saturated capacity rising from 0.69 to
1.98~mg/g as lime proportion increased \parencite{Li2017}.

For Ca-based slag, retention proceeds by calcium phosphate precipitation.
Downstream slag segments were more efficient than upstream segments because they
carried higher $\mathrm{Ca^{2+}}$ concentrations and pH, favouring precipitation,
and phosphorus was shown to be removed in a moving front as $\mathrm{Ca^{2+}}$ and
slag pH-buffer capacity are consumed \parencite{Wang2020d}. Consistent with this,
lower pH improved DRP sorption for Fe-based media, whereas higher/basic pH
improved sorption and precipitation-based retention for Ca-based media
\parencite{Mendes2022}. The pH dependence of Ca systems is sharp in the field:
treated-water pH declined over time and dissolved P removal dropped dramatically
once pH fell below 8.5, with 100\% removal sustained only until cumulative loading
reached about 500~mg~P/kg slag \parencite{Penn2020}. The same buffering principle
underlies subsurface CaCO$_3$ precision treatments, which produced an immediate
rise in discharge-water pH by targeting hydrologically active macropores and
forming a buffering coating near the drains \parencite{Dalhem2019}.

\subsection{Governing Conditions: pH, Coexisting Ions, Temperature, and Particle Size}

Media performance is strongly modulated by water chemistry and grain size.
Across inorganic adsorbents, phosphate adsorption decreased with increasing pH
and acidic conditions generally favoured adsorption, although several materials
retained strength near neutral pH \parencite{Li2020}. For lime-amended pellets
the opposite held---adsorption increased with pH, with equilibrium sorbate rising
from 0.3--1.37~mg/g at pH~4 toward theoretical values at pH~10---reflecting the
precipitation pathway of Ca media \parencite{Li2017}. Coexisting sulfate and
bicarbonate often reduced phosphate removal, whereas chloride and nitrate usually
had little effect for inner-sphere adsorbents \parencite{Li2020}; correspondingly,
nitrate and sulfate caused little interference for steel and mineral media
(variations typically $<10\%$), but humic acid and woodchip dissolved organic
carbon reduced adsorption by 5--47\%, and real drainage waters caused reductions
of 2.3--9.7\% (Arlington site) to 9--21\% (Baltic site) \parencite{Sellner2019b}.
Temperature has a measurable positive effect: phosphate adsorption capacities
increased by factors of 1.2--2.5 as temperature rose from $5^{\circ}$C to
$30^{\circ}$C, consistent with the endothermic nature of the process
\parencite{Sellner2019b}. Smaller particle sizes generally increased DRP sorption
capacity and retention and could prolong media lifetime \parencite{Mendes2022},
a result echoed by the finding that finer particle size and the presence of Fe
oxides (confirmed by XRD) enhanced sorption capacity \parencite{Pugliese2023}
and that furnace slag fine particles achieved the greatest sorption
\parencite{Sanford2016}. Media with varied chemical composition or intra-granular
diffusion (e.g.\ Fe-based CFH-12, Ca-based Filtralite-P) showed enhanced retention
\parencite{Mendes2022}. Successful field use therefore requires attention to
near-neutral pH performance, selectivity, stability, reusability, recovery, and
cost \parencite{Li2020}.

\subsection{Kinetics, Contact Time, and Hydraulic Behaviour}

Phosphate uptake by the better media is rapid, which permits short contact times
in flow-through designs. Column experiments indicated that contact times of only
10--15~s were sufficient for near-complete removal of phosphate from drainage
water spiked with 1~ppm phosphate-P \parencite{Allred2017}, and iron-based
materials reduced soluble P below detection or removed $\geq 90\%$ even at high
flow rates and short contact times \parencite{Allred2014}. Limestone adsorbed
phosphate very rapidly, reaching 98\% of its 24~h capacity within 10~min, while
within 6~h steel slag, small steel chips, and medium steel chips reached 73\%,
67\%, and 51\% of their 24~h capacities \parencite{Sellner2019b}. Lime-free
fly-ash pellets reduced solution P by 80\% in 30~s and absorbed almost all P
within 10~min, whereas lime-bearing pellets required 15--28~min for an 80\%
reduction \parencite{Li2017}. Both granular iron-oxide media and steel-furnace
slag removed SRP below detection for up to five hours of operation; after 24~h the
steel-furnace slag removed about 50\% of SRP while the iron-oxide media still
removed to below detection, and after continuous operation the slag saturated
within seven days whereas the iron-oxide media still removed 20\% SRP after ten
days \parencite{Hauda2020}. Removal generally increased with residence and
empty-bed contact time (EBCT): in slag columns the difference between 50~s and
17~s residence time was significant but that between 17~s and 7~s was not
\parencite{Sanford2016}; longer EBCT improved removal, extended exhaustion time,
and raised modeled capacity for Alcan \parencite{Kumari2025}; and P removal rose
with retention time in staged slag segments \parencite{Wang2020d}. However,
granular iron-oxide media capacity was minimally affected by flow rate and by
influent SRP typical of subsurface drainage, and extending EBCT beyond the tested
range did little to increase capacity \parencite{Hauda2020}. Adequate hydraulic
conductivity underpins field use: all four iron-based materials exceeded the
stormwater sand-filter criterion of $1.0\times10^{-3}$~cm/s
\parencite{Allred2014}, ZVI and SMI materials generally exceeded that threshold
\parencite{Allred2012}, and silica sand, blast-furnace slag, zeolite, CS-AC, and
the CKD/sand blends fell in or near the target range for permeable filter media
\parencite{Agrawal2011}.

\subsection{Desorption, Regeneration, and Media Longevity}

The durability of retention differs sharply between the two families. Fe-based
media exhibit low desorption: phosphate adsorbed or precipitated onto the four
iron-based materials was not readily released, especially for SMI and IOH
\parencite{Allred2014}, and steel chips showed no apparent desorption while
limestone released 9.6\% and steel slag only 0.6\% of adsorbed phosphate over ten
days \parencite{Sellner2019b}. The low desorption potential of steel by-products
was highlighted as beneficial for long-term field performance
\parencite{Sellner2019b}, and reviews note that Fe-based media show negligible
desorption whereas Ca-based media more often exhibit weakly bound DRP prone to
release, particularly at low DRP concentrations or high flow rates
\parencite{Mendes2022}. Several media are regenerable: a NaOH flush plus water
rinse remobilised about 81\% of phosphate adsorbed during a column regeneration
test of SG-IOH \parencite{Allred2017}, and 0.1~M KOH restored Alcan capacity above
93\% over six cycles \parencite{Kumari2025}; no phosphate desorption was detected
from lime-amended pellets after 24~h \parencite{Li2017}. Spent field slag retained
useful capacity---with a 10~min retention time the spent slag still had
considerable P-sorption potential, and laboratory tests showed little
water-extractable P release (seven of nine samples released none; two released
only 0.047 and 0.09~mg/L) \parencite{Shedekar2020}. Longevity is ultimately set by
cumulative loading and the moving-front consumption of $\mathrm{Ca^{2+}}$ and
buffering capacity in slag \parencite{Wang2020d}: dissolved P removal remained at
100\% only until roughly 500~mg~P/kg slag before becoming variable and sometimes
negative \parencite{Penn2020}. Field structures also underperform their design
basis: at equal cumulative loading of 200~mg/kg a field structure removed
50~mg/kg (30\%) of dissolved P versus 125~mg/kg ($\sim$65\%) for the
pre-construction slag used in design \parencite{Shedekar2020}.

\subsection{Field-Scale Phosphorus Removal Structures and Drainage Models}

Field structures demonstrate that reactive media can treat real drainage discharge
while confirming the retention-time and pH limits identified in the laboratory.
A subsurface tile-drain P removal structure removed 19.2~kg of dissolved P,
equal to 55\% of the cumulative dissolved P input, while a companion ditch-style
structure removed 0.9~kg (37\% of inflow load); precipitation-induced events
delivered nearly 97\% of the load while accounting for 65\% of flow volume, and
both field structures underperformed laboratory flow-through experiments, with the
ditch structure showing clogging and non-uniform flow \parencite{Penn2020}. A
ditch-style structure using aluminium-treated steel slag removed about 27\% of
dissolved P and 50\% of total P (2.4~kg and 9.4~kg), lowering outlet flow-weighted
mean concentrations by 32.2\% (dissolved P) and about 41\% (total P); removal was
strong early but became highly variable after about January 2017, though some
later events still reached 100\%, and the Al-treated slag outperformed
non-treated slag at the same site \parencite{Shedekar2020}. A relatively small P
removal structure containing 1700~kg of activated Al in a concrete tank, operated
bottom-up and saturated, treated the entire flow volume without restricting field
drainage and removed 30\% of the 2-yr dissolved P load (240~g) and 29\% of the
cumulative nitrate-N load (8.8~kg); during a 9-day dissolved P spike test it
removed 74\% of dissolved P (770~g), about three times the mass removed during the
2-yr natural-load period \parencite{Penn2024}. In the initial small-scale field
test of SG-IOH-O, phosphate removal averaged 89\% over the first 200~days before
declining to 40\% over the last 80~days, and in concurrent tests SG-IOH-O averaged
75\% removal versus 34\% for the regenerated material; neither formulation
adversely affected water pH \parencite{Allred2017}. Not all field trials succeeded:
a GBFS/CS-AC/CKD/sand/zeolite filter blend in a Texas field did not significantly
reduce $\mathrm{PO_4^{3-}}$ and instead offloaded phosphate above flow rates of
0.037~L~s$^{-1}$, although the same blend removed chlorothalonil (median decrease
69\%, up to 96\%) and otherwise maintained good effluent quality
\parencite{Agrawal2011b}. Flow-through and drainage-model studies provide the
design basis for such structures: staged slag experiments showed discrete P
removal declining with time and loading (single-segment removal averaging 33\% at
2.5~mg~L$^{-1}$ and 24\% at 5.0~mg~L$^{-1}$ inflow), rising with the number of
segments and retention time, and yielded an estimated slag removal capacity of
about 60--61~mg~kg$^{-1}$ at a 30~min retention time, information used to size a
filter unit for a given P load, lifetime, and target efficiency
\parencite{Wang2020d}. A continuous drainage model treating actual farm drainage
water confirmed that clinoptilolite around the drains gave the best performance
among configurations, with a plant (\textit{Phragmites australis}) further raising
removal by 10--20\% \parencite{Sadeghi2021}.

\subsection{Co-contaminant Removal and Integration with Other Practices}

Reactive media can address phosphate together with other drainage contaminants,
enabling stacked treatment trains. In batch tests three SMI materials removed at
least 86\% of nitrate and ZVI1 achieved an 87.8\% atrazine decrease; a two-stage
column of SMI3 followed by ZVI1 removed nitrate, phosphate, and atrazine together,
supporting combined use of selected SMI and ZVI media \parencite{Allred2012}.
Because Fe-based P adsorbents produce negligible acidity or alkalinity effects,
they pair naturally with woodchip bioreactors that remove nitrate: steel
by-product filters were proposed for coupling with bioreactors to remove nitrate
and phosphate simultaneously \parencite{Sellner2019b}. The stacked bioreactor/P
removal design illustrates the hydraulic trade-off: a stacked wood-chip bioreactor
removed 93\% of the nitrate-N and 29\% of the dissolved P entering it, but because
93\% of the 21~ML total tile discharge bypassed the bioreactor untreated this
translated to only a 5.4\% nitrate-N load reduction and 15~g of dissolved P
removed, whereas the P removal structure treating the full flow removed nearly
ten times more nitrate-N \parencite{Penn2024}. Complementary edge-of-field
practices target nitrate directly: all six saturated riparian buffers were
effective at removing nitrate from tile outlet water, with annual loads removed of
13--179~kg~N over drainage areas of 3.4--40.5~ha, removal effectiveness of 8--84\%,
and removal rates of 0.27--4.13~g~N~m$^{-1}$~d$^{-1}$ \parencite{Jaynes2019}.

\subsection{Cost-effectiveness and Practical Implementation}

Cost analyses indicate that reactive-media treatment is feasible and can be
optimised by media choice and end-of-life strategy. For a subsurface drainage
case study, the total media volume required to remove 1.66~kg SRP was 4.22~m$^3$
for granular iron-oxide media versus 15.5~m$^3$ for steel-furnace slag, and annual
cost per hectare ranged from \$1641 (scenario~A) to \$906 (scenario~C), with the
lowest-cost option being single-use steel-furnace slag disposed after use; onsite
regeneration of the iron-oxide media was slightly more expensive but the
methodology transfers to other sites \parencite{Hauda2020}. Saturated riparian
buffers, a complementary practice, had a mean equalised annual cost of
US\$213.83, or US\$2.94~kg$^{-1}$~N removed, and were described as comparatively
simple, cheap, and quick to install with little maintenance
\parencite{Jaynes2019}. Multiple studies conclude that low-cost by-products are
practical media: recycled steel by-products were identified as highly efficient
phosphate adsorbents across a wide range of temperature, pH, organic matter, and
real drainage water conditions and are recommended for subsurface drainage
\parencite{Sellner2019b}; furnace slag small particles ($\sim$20--70\% P removal)
may suit an inline tile-drainage treatment system \parencite{Sanford2016}; and
durable fly-ash pellets could serve as a tile-outlet filter for drainage-water
treatment \parencite{Li2017}. Fe-rich materials are recommended for effective
field-scale application because they provide extended, relatively high P retention
with lower sensitivity to variable flow and P concentration
\parencite{Pugliese2023}, and the readily acquired constituents of slag/zeolite
blends make them attractive for large-scale implementation despite the phosphate
offloading observed for one blend \parencite{Agrawal2011b}.

\subsection{Complementary Drainage-Water Chemistry and Monitoring Context}

Related work situates reactive-media P removal within broader drainage-water
management and monitoring. Subsurface CaCO$_3$ precision treatments in acid
sulfate soils produced durable improvements in discharge-water chemistry: the
large CaCO$_3$ dose kept pH $+0.4$ units above reference four years after
application, aluminium decreased by about 95\% (to $\sim$1~mg~L$^{-1}$), and Fe
speciation shifted toward schwertmannite and amorphous iron oxides under higher
pH and lower ORP; Ca(OH)$_2$ treatments raised pH and cut Al by up to 99\%
immediately but were short-lived, and sulfur was unaffected \parencite{Dalhem2019}.
CaCO$_3$ amendment also influenced soil water behaviour, with high CaCO$_3$
helping reduce swelling and excess field-capacity water retention irrespective of
EC \parencite{He2015}. Where drainage doubles as subirrigation supply, a new
performance method reported water-application efficiencies and uniformity metrics
(e.g.\ distribution uniformity of the low quarter of 48--78\%) that improved with
longer irrigation durations or slower flow rates \parencite{Jia2017}. Subsurface
pipe drainage markedly altered soil ion contents---lowering $\mathrm{Na^+}$,
$\mathrm{Ca^{2+}}$, and $\mathrm{Mg^{2+}}$ while $\mathrm{SO_4^{2-}}$ became the
most abundant ion---and HYDRUS-2D/UnsatChem reproduced ion discharge with
calibration $R^2 > 0.87$, showing that ion-composition-based drainage designs
outperform total-salt designs \parencite{Liu2025b}. For monitoring, field
electrical conductivity was a precise, cost-effective estimator of the cumulative
concentration of major ions, with $\mathrm{HCO_3^-}$, $\mathrm{Ca^{2+}}$, and
$\mathrm{Mg^{2+}}$ accounting for 90\% of calculated conductivity
\parencite{Zimmerman2017}. Finally, electrochemical treatment of drainage brine
was optimised by response-surface methods to $3.6~\mathrm{mol\,L^{-1}}$ brine,
pH~9, 12.2~V, and 22.5~cm electrode immersion depth, substantially improving
performance \parencite{Yang2023}.